\chapter{Źródła i weryfikacja}\markboth{Źródła i weryfikacja}{Źródła i weryfikacja}

Ta książka ma niewiele twierdzeń o świecie, ale każde, które ma, powinno dać się sprawdzić.
Ten dodatek opisuje, jak to robić, i wyjaśnia, czego znaczek [?] w tekście nie znaczy.

\section{Co znaczy znaczek [?]}
Pomarańczowy znaczek [?] stoi przy twierdzeniu, które dotyczy świata poza tekstem, czyli zdrowia,
pamięci, czasu czytania innych ludzi, prawa albo liczby spoza repozytorium, i którego autor nie
sprawdził w pierwotnym źródle. Znaczek nie oznacza, że twierdzenie jest błędne. Oznacza, że nie
wiemy, czy jest prawdziwe, bo nikt tego nie zapisał w sposób, który dałoby się sprawdzić.

Twierdzenie usuwa się z listy na dwa sposoby. Pierwszy: zamienia się je na arytmetykę pokazaną
w tekście albo na liczbę policzoną z danych, czyli na coś, co czytelnik może odtworzyć. Drugi:
podaje się źródło z wystarczającymi danymi, by je odnaleźć, i zapisuje się, co w nim faktycznie
stwierdzono. Usunięcie znaczka bez żadnej z tych czynności oznacza, że twierdzenie wciąż jest
niesprawdzone, tylko przestało być widoczne.

W tej wersji nie podaję tytułów ani lat źródeł, których nie sprawdzono w tej sesji. Zamiast
nazwisk czytelnik znajdzie tu typy źródeł, których należy szukać.

\section{Typy źródeł}
\begin{enumerate}
\item \textbf{Badania pierwotne} z eksperymentami porównującymi różne sposoby uczenia się lub
czytania z kontrolą. Najpierw sprawdzaj te, bo tylko one mogą odpowiedzieć na pytanie
przyczynowe.
\item \textbf{Przeglądy systematyczne i metaanalizy}, które zbierają wiele badań. Ich wartość
zależy od tego, jak wybrano badania i czy autorzy uwzględnili te negatywne.
\item \textbf{Dane statystyczne} instytucji, które mierzą czytanie i rynek wydawniczy. Dobre do
pytań o liczby książek i ich długość; sprawdź, jak zdefiniowano pojęcie \glqq książka\grqq{}.
\item \textbf{Pomiary własne} czasu z zegarem i liczby słów. Najlepsze do twojego tempa, ale
zawsze z jednym zastrzeżeniem: mierzysz jeden tekst, nie wszystkie.
\item \textbf{Relacje osobiste i poradniki.} Przydatne jako hipotezy do sprawdzenia, najsłabsze
jako dowód.
\end{enumerate}

\section{Jak sprawdzać twierdzenie o czytaniu lub pamięci}
Zacznij od pracy pierwotnej, a nie od omówienia, bo omówienia zwykle upraszczają wynik. Potem
zapytaj, czy porównanie miało grupę kontrolną, ile osób badano i kogo. Zwróć uwagę, co dokładnie
zmierzono: wynik testu, czy samoocenę. Sprawdź, czy wynik był powtórzony w innym badaniu, i jak duży
był efekt, w jakich jednostkach i względem czego. Dopiero wtedy decyduj, czy twierdzenie wolno
przenieść do swojego planu.

\section{Lista kontrolna}
Przed użyciem jakiegokolwiek twierdzenia o czytaniu, pamięci lub czasie przejdź przez osiem punktów:
\begin{enumerate}
\item Czy znasz źródło pierwotne, a nie tylko omówienie?
\item Czy porównanie ma grupę kontrolną?
\item Czy podano wielkość próby i to, kogo badano?
\item Czy wiesz, co zmierzono: wynik czy poczucie?
\item Czy wynik powtórzono w innym badaniu?
\item Czy znasz wielkość efektu i jednostki, w których ją podano?
\item Czy umiesz przeliczyć twierdzenie na własne liczby, na przykład z dodatku B?
\item Czy zapisałeś odpowiedź i źródło w swoim dzienniczku, żeby nie szukać ich ponownie?
\end{enumerate}
Jeśli odpowiedź na którekolwiek z pytań brzmi \glqq nie\grqq{}, traktuj twierdzenie jako hipotezę,
a nie jako podstawę decyzji o czasie, który masz.
